%=====================================================================
\chapter{Problem Formulation and Uninformed Search}\label{ch:uninformed}
%=====================================================================
\locator{2}{Part II --- Problem Solving by Search}

\begin{outcomes}
By the end of this chapter you will be able to:
\begin{tight}
  \item \textbf{Formulate} an informal task as a search problem, choosing a state
        representation and \textbf{justifying} what it leaves out.
  \item \textbf{Distinguish} the state space from the search tree, and
        \textbf{explain} why the tree can be infinite when the space is finite.
  \item \textbf{Trace} breadth-first, depth-first, depth-limited, iterative
        deepening and uniform-cost search by hand on a small graph.
  \item \textbf{Compare} the five strategies on completeness, optimality, time and
        space, and \textbf{choose} one for a given problem with a reason.
  \item \textbf{Explain} why uninformed search cannot scale, and \textbf{quantify}
        the growth that defeats it.
\end{tight}
\end{outcomes}

\begin{prereq}
\chref{python} (the \texttt{Problem} interface, nodes, the three frontiers,
\texttt{frozenset} states) and \chref{agents} (the goal-based agent, whose
deliberation is exactly the search of this chapter). Every algorithm here is the
same twelve-line body with a different frontier.
\end{prereq}

\begin{keyterms}
problem formulation $\bullet$ state space $\bullet$ search tree $\bullet$ node
expansion $\bullet$ frontier $\bullet$ explored set $\bullet$ redundant path
$\bullet$ branching factor $\bullet$ depth $\bullet$ completeness $\bullet$
optimality $\bullet$ breadth-first search $\bullet$ depth-first search $\bullet$
depth-limited search $\bullet$ iterative deepening $\bullet$ uniform-cost search
$\bullet$ combinatorial explosion
\end{keyterms}

\section{From a Task to a State Space}

A delivery manager says: \emph{get the release out, and don't start anything before
the things it depends on are finished.} That is a task, not a problem. Turning it
into a problem means deciding five things --- the five of \chref{python} --- and
the decisions are not mechanical.

\begin{examplebox}[Formulating the release-ordering task]
\begin{tight}
  \item \emph{States.} The set of finished tasks. A \texttt{frozenset}, so that two
        routes that finished the same work collapse to one state.
  \item \emph{Initial state.} The empty set.
  \item \emph{Actions.} In state $s$, the tasks not yet done whose dependencies are
        all in $s$ --- the \texttt{ready} function of \chref{python}.
  \item \emph{Result.} $s \cup \{t\}$. A new set; the old one is never touched.
  \item \emph{Goal test.} Is every task in $s$?
  \item \emph{Step cost.} The task's duration in days.
\end{tight}
\end{examplebox}

Notice what the formulation \emph{throws away}: who does each task, the calendar,
holidays, the fact that two tasks by the same person cannot overlap. Each omission
is a modelling decision that makes the problem solvable and the answer less
directly useful. \chref{csp} puts the people back in, and it costs a different
algorithm to do it.

\begin{conceptbox}
The single most valuable habit in this part of the course: \emph{write the
formulation down before writing the search}, and beside each state component write
why the goal test needs it. Anything the goal test cannot see is a candidate for
deletion, and every deletion shrinks the space multiplicatively.

Beginners reach for a cleverer algorithm when a search is too slow. Professionals
reach for the formulation first, because that is where the factors of a hundred
are. \chref{python}'s worked example measured one such factor; this chapter's lab
measures another.
\end{conceptbox}

\section{The State Space and the Search Tree}

These are two different objects and confusing them causes both of the classic bugs.

The \term{state space} is a graph: nodes are states, edges are actions. It is a
property of the problem. For this project it has $22$ reachable states.

The \term{search tree} is what the algorithm builds as it explores. Its root is the
initial state, and its branches are paths. The same state appears in it once for
every distinct path that reaches it --- so the tree has $184$ nodes where the graph
has $22$. \dsref{space-vs-tree} shows both for a three-task fragment.

\dsfig{%
\begin{tikzpicture}[font=\scriptsize,
  st/.style={circle, draw=primary, line width=0.9pt, fill=primary!8,
             text=primary!55!black, font=\tiny, inner sep=1.4pt,
             minimum size=0.95cm, align=center},
  tn/.style={circle, draw=secondary, line width=0.9pt, fill=secondary!8,
             text=secondary!50!black, font=\tiny, inner sep=1.4pt,
             minimum size=0.95cm, align=center},
  dup/.style={circle, draw=accent, line width=1.2pt, fill=accent!12,
              text=accent!55!black, font=\tiny, inner sep=1.4pt,
              minimum size=0.95cm, align=center},
  e/.style={-{Stealth[length=1.6mm]}, primary!50, line width=0.8pt},
  f/.style={-{Stealth[length=1.6mm]}, secondary!50, line width=0.8pt},
  lb/.style={font=\tiny, text=neutral, inner sep=1pt}]

% ---- the state space: a graph
\node[font=\scriptsize\bfseries, text=primary] at (-4.3,1.5)
  {STATE SPACE (a graph)};
\node[st] (e0) at (-6.0,0.4)  {$\varnothing$};
\node[st] (a)  at (-4.4,1.0)  {a};
\node[st] (b)  at (-4.4,-0.2) {b};
\node[st] (ab) at (-2.7,0.4)  {ab};
\node[st] (abc) at (-1.1,0.4) {abc};
\draw[e] (e0) -- node[lb, above left] {a} (a);
\draw[e] (e0) -- node[lb, below left] {b} (b);
\draw[e] (a)  -- node[lb, above right] {b} (ab);
\draw[e] (b)  -- node[lb, below right] {a} (ab);
\draw[e] (ab) -- node[lb, above] {c} (abc);
\node[font=\tiny\itshape, text=primary, align=center] at (-3.5,-1.4)
  {5 states, 5 edges.\\ Two ways into \texttt{ab}; it is \emph{one} state.};

% ---- the search tree
\node[font=\scriptsize\bfseries, text=secondary] at (3.1,1.5)
  {SEARCH TREE (paths)};
\node[tn] (r)   at (3.1,0.75)  {$\varnothing$};
\node[tn] (ta)  at (1.7,-0.35) {a};
\node[tn] (tb)  at (4.5,-0.35) {b};
\node[dup] (tab1) at (1.7,-1.6) {ab};
\node[dup] (tab2) at (4.5,-1.6) {ab};
\node[tn] (tc1) at (1.7,-2.85) {abc};
\node[tn] (tc2) at (4.5,-2.85) {abc};
\draw[f] (r) -- node[lb, above left] {a} (ta);
\draw[f] (r) -- node[lb, above right] {b} (tb);
\draw[f] (ta) -- node[lb, left] {b} (tab1);
\draw[f] (tb) -- node[lb, right] {a} (tab2);
\draw[f] (tab1) -- node[lb, left] {c} (tc1);
\draw[f] (tab2) -- node[lb, right] {c} (tc2);

\draw[accent, line width=1pt, dash pattern=on 2pt off 1.6pt]
  (tab1) to[bend right=32] (tab2);
\node[font=\tiny\bfseries, text=accent, align=center] at (3.1,-1.98)
  {same state, twice};

\node[rounded corners=3pt, fill=accent!8, draw=accent, line width=0.9pt,
      text=accent!55!black, font=\scriptsize, text width=13.4cm,
      align=flush left, inner sep=6pt] at (-0.4,-4.3)
  {\textbf{The explored set is what turns the tree back into the graph.} Keep a
   set of states already generated and refuse to add one twice, and the search
   visits $5$ nodes instead of $7$ --- here a small saving, on the real project the
   difference between $22$ and $184$. Omit it and a state space with a cycle gives
   an \emph{infinite} search tree, which is the first classic bug.};
\end{tikzpicture}}%
{The state space and the search tree for a three-task fragment where \texttt{c}
depends on both \texttt{a} and \texttt{b}. The tree repeats every state once per
path that reaches it.}{space-vs-tree}

\begin{alertbox}[The two classic bugs, and their symptoms]
\textbf{No explored set.} Symptom: the search never terminates, or returns a path
that revisits a state. Cause: the search tree of a cyclic space is infinite. Fix:
a \texttt{set} of states, tested before pushing.

\smallskip
\textbf{Goal-testing on generation instead of on expansion.} Symptom: breadth-first
search returns a correct shallowest path (fine), but uniform-cost search returns a
\emph{sub-optimal} one. Cause: when a node is generated, a cheaper path to the same
state may still be in the frontier. Breadth-first search may test on generation
because all edges cost the same; uniform-cost search must test when the node is
\emph{popped}. This is the subtlest distinction in the chapter and
Section~\ref{sec:ucs} returns to it.
\end{alertbox}

\section{Breadth-First Search}

Expand the shallowest unexpanded node: a first-in first-out frontier, so
\texttt{deque.popleft()}.

Breadth-first search is \term{complete} --- if a solution exists at finite depth it
will be found --- and it returns the \emph{shallowest} solution. That is an optimal
solution only when every action costs the same. With durations of $1$ to $5$ days,
the shallowest plan is not the cheapest, and it is a mistake to call breadth-first
search optimal without that qualification.

Its cost is the problem. With branching factor $b$ and solution depth $d$, it
generates $O(b^{d})$ nodes and --- fatally --- must \emph{store} $O(b^{d})$ of them,
because the whole frontier is one level wide.

\section{Depth-First and Depth-Limited Search}

Expand the deepest unexpanded node: a last-in first-out frontier, so
\texttt{list.pop()}, or simple recursion.

The gain is memory: only the current path and its siblings need to be held, which
is $O(bm)$ for maximum depth $m$ rather than $O(b^{d})$. The loss is both
completeness and optimality --- in a space with cycles and no explored set,
depth-first search follows one branch forever.

\term{Depth-limited search} fixes the infinite descent by refusing to go below
depth $\ell$. Now it terminates, but it is incomplete in a new way: if the shallowest
solution is deeper than $\ell$ it reports failure on a solvable problem, which is
the worst possible answer.

\section{Iterative Deepening}

Run depth-limited search with $\ell = 0$, then $1$, then $2$, and so on until a
solution is found. This sounds wasteful and is the standard answer for large
uninformed problems, because it has breadth-first search's completeness and
shallowest-solution guarantee with depth-first search's memory.

The apparent waste is small. The work is dominated by the deepest level, because
the levels grow geometrically:
\[
  N(\text{IDS}) = (d+1)b^{0} + d\,b^{1} + (d-1)b^{2} + \cdots + 1\cdot b^{d}
\]
and for $b = 10$, $d = 5$ this is $123{,}456$ nodes against breadth-first search's
$111{,}111$ --- about 11\% more work, in exchange for storing $50$ nodes instead of
$100{,}000$. \dsref{ids-growth} shows why the ratio improves as $b$ grows.

\dsfig{%
\begin{tikzpicture}
\begin{axis}[
  width=11.6cm, height=6.0cm,
  xlabel={branching factor $b$}, ylabel={nodes generated},
  xlabel style={font=\scriptsize}, ylabel style={font=\scriptsize},
  tick label style={font=\tiny},
  ymode=log, log basis y=10,
  xmin=2, xmax=12, ymin=10, ymax=4e6,
  grid=both, major grid style={gray!18}, minor grid style={gray!8},
  legend style={font=\tiny, at={(0.03,0.97)}, anchor=north west,
                draw=gray!40, fill=white},
  legend cell align=left]

% breadth-first: (b^(d+1)-1)/(b-1) for d=5, i.e. nodes to depth 5
\addplot[primary, line width=1.2pt, mark=*, mark size=1.3pt] coordinates {
  (2,63) (3,364) (4,1365) (5,3906) (6,9331) (7,19608) (8,37449)
  (9,66430) (10,111111) (11,177156) (12,271453)};
\addlegendentry{breadth-first, nodes stored $\approx b^{d}$}

% iterative deepening total work, d=5
\addplot[greenhl, line width=1.2pt, mark=square*, mark size=1.3pt,
         dash pattern=on 3pt off 2pt] coordinates {
  (2,120) (3,543) (4,1818) (5,4881) (6,11196) (7,22875) (8,42798)
  (9,74733) (10,123456) (11,194871) (12,296130)};
\addlegendentry{iterative deepening, total work}

% memory: b*d
\addplot[accent, line width=1.2pt, mark=triangle*, mark size=1.6pt] coordinates {
  (2,10) (3,15) (4,20) (5,25) (6,30) (7,35) (8,40) (9,45) (10,50)
  (11,55) (12,60)};
\addlegendentry{iterative deepening, nodes stored $= bd$}
\end{axis}
\end{tikzpicture}}%
{Iterative deepening against breadth-first search at solution depth $d=5$. The
repeated work is a modest overhead that shrinks as the tree widens --- 91\% extra
at $b=2$, 11\% at $b=10$, 9\% at $b=12$ --- while the memory requirement falls from
exponential to linear.}{ids-growth}

\section{Uniform-Cost Search}
\label{sec:ucs}

Expand the node with the lowest \emph{path cost} $g(n)$, not the shallowest: a
priority-queue frontier, so \texttt{heapq}.

This is the first algorithm in the book that is genuinely optimal in cost, and it
is optimal for a reason worth stating precisely: when a node is \emph{popped}, every
remaining path to it costs at least as much, because the frontier is ordered by
cost and costs are non-negative. So the first time the goal is popped, it is popped
by its cheapest path.

That argument is exactly why the goal test must happen on \emph{popping} and not on
generating, and why every edge cost must be non-negative. Break either condition
and uniform-cost search quietly returns sub-optimal answers.

Uniform-cost search is breadth-first search when all costs are equal, and it is
Dijkstra's algorithm. In \chref{informed} it becomes $A^{*}$ by adding one term.

\smallskip
\noindent\begin{longtable}{@{}L{2.6cm}L{2.05cm}L{1.95cm}L{2.5cm}L{1.85cm}L{1.85cm}@{}}
\toprule
\textbf{Strategy} & \textbf{Complete?} & \textbf{Cheapest?} & \textbf{Time} &
\textbf{Space} & \textbf{Frontier} \\
\midrule
\endfirsthead
\toprule
\textbf{Strategy} & \textbf{Complete?} & \textbf{Cheapest?} & \textbf{Time} &
\textbf{Space} & \textbf{Frontier} \\
\midrule
\endhead
Breadth-first & Yes & Only if costs equal & $O(b^{d})$ & $O(b^{d})$ &
\texttt{deque} \\
Depth-first & No (cycles) & No & $O(b^{m})$ & $O(bm)$ & \texttt{list} \\
Depth-limited & No (if $\ell<d$) & No & $O(b^{\ell})$ & $O(b\ell)$ &
\texttt{list} \\
Iterative deepening & Yes & Only if costs equal & $O(b^{d})$ & $O(bd)$ &
\texttt{list} \\
Uniform-cost & Yes & \textbf{Yes} & $O(b^{1+C^{*}/\epsilon})$ &
same as time & \texttt{heapq} \\
\bottomrule
\caption{The five uninformed strategies. $b$ is the branching factor, $d$ the
solution depth, $m$ the maximum depth, $C^{*}$ the optimal cost and $\epsilon$ the
smallest edge cost. Note that only uniform-cost search is optimal in cost.}
\label{tab:uninformed}
\end{longtable}

\section{Why Uninformed Search Cannot Scale}

Every entry in the time column of Table~\ref{tab:uninformed} is exponential, and no
choice of frontier changes that. The reason is arithmetic: at branching factor $10$,
depth $10$ is $10^{10}$ nodes. At a million nodes per second that is about three
hours, and depth $12$ is twelve days.

This is the \term{combinatorial explosion} that ended the first wave of AI research
(\chref{introduction}), and it is not an implementation problem. These algorithms
are \emph{uninformed}: they know nothing about the problem except how to generate
successors and recognise a goal, so they cannot prefer a promising branch. The only
escape is to give the search a way to estimate how far a state is from the goal,
which is \chref{informed}.

\begin{worked}[Tracing all four strategies on one small graph]
The three-task fragment of \dsref{space-vs-tree}, with costs: \texttt{a} costs $3$,
\texttt{b} costs $1$, \texttt{c} costs $2$, and \texttt{c} requires both. Goal:
\texttt{abc}. Expansion order, with the frontier shown after each step.

\smallskip
\textbf{Breadth-first} (first-in first-out, explored set on).
\begin{tight}
  \item Pop $\varnothing$. Push \texttt{a}, \texttt{b}. Frontier:
        [\texttt{a}, \texttt{b}].
  \item Pop \texttt{a}. Push \texttt{ab}. Frontier: [\texttt{b}, \texttt{ab}].
  \item Pop \texttt{b}. Its only successor is \texttt{ab}, already in the explored
        set --- \emph{not} pushed. Frontier: [\texttt{ab}].
  \item Pop \texttt{ab}. Push \texttt{abc}. Frontier: [\texttt{abc}].
  \item Pop \texttt{abc}. Goal. Path \texttt{a},\texttt{b},\texttt{c}, depth $3$,
        cost $3+1+2=6$.
\end{tight}
Nodes expanded: $4$. Without the explored set, step~3 would have pushed a second
\texttt{ab} and everything below it would be built twice.

\smallskip
\textbf{Depth-first} (last-in first-out).
Pop $\varnothing$, push \texttt{a}, \texttt{b}; pop \texttt{b} (the newest); push
\texttt{ab}; pop \texttt{ab}; push \texttt{abc}; pop \texttt{abc}. Goal in $4$
expansions, path \texttt{b},\texttt{a},\texttt{c}. Same cost here by coincidence;
on the full project it is worse.

\smallskip
\textbf{Iterative deepening.}
$\ell=0$: expand $\varnothing$, fail ($1$ node). $\ell=1$: $\varnothing$,
\texttt{a}, \texttt{b}, fail ($3$). $\ell=2$: $\varnothing$, \texttt{a},
\texttt{ab}, \texttt{b}, fail ($4$). $\ell=3$: reaches \texttt{abc}. Total about
$12$ node generations against breadth-first search's $5$ --- but the deepest
frontier ever held is $2$ nodes, not $2$ levels.

\smallskip
\textbf{Uniform-cost} (cheapest first). Frontier as
$(\text{cost},\text{state})$.
\begin{tight}
  \item Pop $(0,\varnothing)$. Push $(3,\texttt{a})$, $(1,\texttt{b})$.
  \item Pop $(1,\texttt{b})$ --- cheaper than \texttt{a}, so the order differs from
        breadth-first search. Push $(4,\texttt{ab})$.
  \item Pop $(3,\texttt{a})$. Successor \texttt{ab} at cost $3+1=4$, no better than
        the $4$ already in the frontier, so nothing improves.
  \item Pop $(4,\texttt{ab})$. Push $(6,\texttt{abc})$.
  \item Pop $(6,\texttt{abc})$. Goal, cost $6$, and now \emph{provably} cheapest.
\end{tight}

\smallskip
\textbf{The point of the trace.} All four found a solution; only uniform-cost
search knows its answer is the cheapest. And notice step~3: had the goal been
tested when \texttt{ab} was \emph{generated} rather than when it was popped, the
search would have committed to whichever path reached it first, which is the bug
the alertbox warns about.
\end{worked}

\begin{pitfall}
\begin{tight}
  \item \textbf{Omitting the explored set.} An infinite search tree on a finite
        space. The symptom is a hang, and the cause is two lines away.
  \item \textbf{Calling breadth-first search optimal.} It returns the
        \emph{shallowest} solution. That is the cheapest only when every action
        costs the same --- which, with task durations, it does not.
  \item \textbf{Goal-testing on generation in uniform-cost search.} Silently
        sub-optimal answers, and the tests pass because the answer is still
        \emph{a} solution.
  \item \textbf{Negative step costs.} Uniform-cost search's optimality proof needs
        non-negative costs. ``Refactoring saves two days'' modelled as a negative
        cost breaks the algorithm, not just the estimate.
  \item \textbf{Depth-limited search reporting failure as impossibility.} It means
        ``no solution within $\ell$''. Treating that as ``no solution'' is a
        reporting bug with real consequences.
  \item \textbf{Believing a faster language will fix $b^{d}$.} A hundredfold speedup
        buys two more levels of depth at $b=10$. Change the formulation or add a
        heuristic.
\end{tight}
\end{pitfall}

%---------------------------------------------------------------------
\begin{lab}[Task Ordering Four Ways]
One problem, four strategies, one table. The search body is the one from
\chref{python}; only the frontier changes. The lab also measures the cost of
dropping the explored set.
\begin{lstlisting}[language=Python]
"""Chapter 4 lab -- uninformed search on the project graph.

Breadth-first, depth-first, iterative deepening and uniform-cost over
the same task-ordering problem. The search body never changes: the
frontier decides the strategy, exactly as Chapter 3 claimed.
"""

import heapq
import itertools
from collections import deque, namedtuple

# duration in days, and the tasks that must finish first
TASKS = {
    "spec":    (3, set()),
    "schema":  (2, {"spec"}),
    "api":     (5, {"schema"}),
    "ui":      (4, {"spec"}),
    "auth":    (3, {"schema"}),
    "tests":   (4, {"api", "auth"}),
    "docs":    (2, {"api"}),
    "uat":     (3, {"tests", "ui"}),
    "release": (1, {"uat", "docs"}),
}
ALL = frozenset(TASKS)
Node = namedtuple("Node", "state parent action cost depth")


def actions(state):
    for t, (_, deps) in TASKS.items():
        if t not in state and deps <= state:
            yield t


def result(state, action):
    return state | {action}


def cost_of(action):
    return TASKS[action][0]


def path(node):
    out = []
    while node.parent is not None:
        out.append(node.action)
        node = node.parent
    return list(reversed(out))


def graph_search(pop_which, explored=True):
    """One body. pop_which selects the strategy.

    'fifo'  -> breadth-first      'lifo' -> depth-first
    'cheap' -> uniform-cost
    """
    start = Node(frozenset(), None, None, 0, 0)
    tie = itertools.count()
    if pop_which == "cheap":
        frontier = [(0, next(tie), start)]
    else:
        frontier = deque([start])
    seen = {start.state} if explored else None
    expanded = 0
    peak = 1
    while frontier:
        if pop_which == "cheap":
            node = heapq.heappop(frontier)[2]
        elif pop_which == "fifo":
            node = frontier.popleft()
        else:
            node = frontier.pop()

        if node.state == ALL:           # goal test on POPPING
            return path(node), node.cost, expanded, peak
        expanded += 1
        for a in actions(node.state):
            nxt = result(node.state, a)
            if explored:
                if nxt in seen:
                    continue
                seen.add(nxt)
            child = Node(nxt, node, a, node.cost + cost_of(a),
                         node.depth + 1)
            if pop_which == "cheap":
                heapq.heappush(frontier, (child.cost, next(tie), child))
            else:
                frontier.append(child)
        peak = max(peak, len(frontier))
    return None, None, expanded, peak


def depth_limited(node, limit, seen, stats):
    """Recursive depth-limited search. Returns a node or None."""
    if node.state == ALL:
        return node
    if node.depth >= limit:
        return None
    stats["expanded"] += 1
    for a in actions(node.state):
        nxt = result(node.state, a)
        if nxt in seen:
            continue
        seen.add(nxt)
        child = Node(nxt, node, a, node.cost + cost_of(a),
                     node.depth + 1)
        found = depth_limited(child, limit, seen, stats)
        if found is not None:
            return found
        seen.discard(nxt)          # allow other branches to reuse it
    return None


def iterative_deepening():
    stats = {"expanded": 0}
    for limit in range(1, len(TASKS) + 2):
        start = Node(frozenset(), None, None, 0, 0)
        found = depth_limited(start, limit, {start.state}, stats)
        if found is not None:
            return path(found), found.cost, stats["expanded"], limit
    return None, None, stats["expanded"], None


def legal(plan):
    done = set()
    for t in plan:
        if not TASKS[t][1] <= done:
            return False
        done.add(t)
    return done == set(TASKS)


if __name__ == "__main__":
    print("tasks: %d   total duration: %d days"
          % (len(TASKS), sum(d for d, _ in TASKS.values())))
    print()
    print("strategy             len  days  expanded  peak")
    print("-" * 49)
    rows = {}
    for name, key in (("breadth-first ", "fifo"),
                      ("depth-first   ", "lifo"),
                      ("uniform-cost  ", "cheap")):
        plan, cost, exp, peak = graph_search(key)
        rows[name.strip()] = (plan, cost, exp, peak)
        assert legal(plan), name
        print("%s %6d %5d %9d %5d"
              % (name, len(plan), cost, exp, peak))
    plan, cost, exp, depth = iterative_deepening()
    assert legal(plan)
    print("%s %6d %5d %9d %5s"
          % ("iter. deepening", len(plan), cost, exp, "-"))

    # every task must be done, so every complete plan costs the same
    costs = {c for (_, c, _, _) in rows.values()}
    assert costs == {sum(d for d, _ in TASKS.values())}
    print()
    print("Every plan costs %d days: all nine tasks are mandatory, so"
          % costs.pop())
    print("cost cannot discriminate. Chapter 5 gives it something to")
    print("discriminate on.")

    # what the explored set is worth
    _, _, with_set, _ = graph_search("fifo", explored=True)
    _, _, without, _ = graph_search("fifo", explored=False)
    print()
    print("breadth-first expansions with explored set:    %d" % with_set)
    print("breadth-first expansions without explored set: %d" % without)
    print("factor: %.1fx more work" % (without / with_set))
    assert without > with_set
\end{lstlisting}
Read the output against Table~\ref{tab:uninformed}: depth-first search holds the
smallest frontier and breadth-first search the largest, every complete plan costs
the same $27$ days, and removing the explored set multiplies the work. The last
number is the one to remember --- it is the whole argument of
\dsref{space-vs-tree}, measured.
\end{lab}

%---------------------------------------------------------------------
\begin{checkpoint}
\begin{tightnum}
  \item A state space has $22$ states and its search tree has $184$ nodes. Explain
        in one sentence where the other $162$ come from, and what single data
        structure removes them.
  \item Why must uniform-cost search test for the goal when a node is popped
        rather than when it is generated? Give the smallest graph on which the
        difference shows.
  \item Iterative deepening re-expands the top of the tree on every iteration. Why
        is the total work still the same order as breadth-first search, and what is
        bought with the repetition?
\end{tightnum}
\end{checkpoint}

%---------------------------------------------------------------------
\begin{chaptersummary}
\begin{tight}
  \item Formulating a problem means choosing states, initial state, actions,
        result, goal test and step cost --- and deciding what to leave out. The
        omissions are where the factors of a hundred live.
  \item The \term{state space} is a graph and belongs to the problem; the
        \term{search tree} is built by the algorithm and repeats each state once
        per path. An \term{explored set} turns the tree back into the graph and is
        what makes search on a cyclic space terminate.
  \item Breadth-first search is complete and returns the shallowest solution ---
        cheapest only when all actions cost the same. It needs $O(b^{d})$ memory,
        which is what defeats it.
  \item Depth-first search needs only $O(bm)$ memory but is neither complete nor
        optimal. Depth-limited search terminates but reports failure on solvable
        problems.
  \item Iterative deepening has breadth-first search's guarantees with depth-first
        search's memory, for about 11\% more work at $b=10$, $d=5$.
  \item Uniform-cost search expands the cheapest path first and is genuinely
        cost-optimal, provided the goal test happens on popping and every step cost
        is non-negative. It is Dijkstra's algorithm, and \chref{informed} turns it
        into $A^{*}$.
  \item All five are exponential in time, because none of them knows anything about
        the problem beyond generating successors. That is what ``uninformed''
        means, and why the next chapter exists.
\end{tight}
\end{chaptersummary}

\begin{reviewq}
  \item \textbf{(MCQ)} Which strategy is guaranteed to return the cheapest
        solution when action costs differ? (a)~breadth-first (b)~depth-first
        (c)~iterative deepening (d)~uniform-cost.
  \item \textbf{(MCQ)} Depth-first search's space complexity is:
        (a)~$O(b^{d})$ (b)~$O(bm)$ (c)~$O(d)$ (d)~$O(b^{m})$.
  \item \textbf{(MCQ)} Omitting the explored set on a state space with cycles makes
        the search: (a)~slower but still correct (b)~non-optimal only
        (c)~potentially non-terminating (d)~unable to find shallow solutions.
  \item \textbf{(MCQ)} Uniform-cost search is equivalent to breadth-first search
        when: (a)~the graph is a tree (b)~all step costs are equal (c)~the
        heuristic is zero (d)~the goal is unique.
  \item \textbf{(Short)} State the difference between the state space and the
        search tree, and give the size of each for the three-task fragment of
        \dsref{space-vs-tree}.
  \item \textbf{(Short)} Explain why iterative deepening is preferred to
        breadth-first search on large problems despite doing more work.
  \item \textbf{(Short)} Why does uniform-cost search require non-negative step
        costs? Describe what goes wrong with a negative edge.
  \item \textbf{(Applied)} On the project graph, add a task \texttt{perf} of
        duration $2$ depending on \texttt{api} and \texttt{ui}, required before
        \texttt{release}. By hand, give the depth of the shallowest complete plan
        and the total cost of every complete plan. Then say which of the five
        strategies would now be the right choice, and why cost-based search is
        still not useful on this problem.
\end{reviewq}
